\section{Validating the Loser-Side Ranking}
\label{sec:validation}

This section validates the award-layer ranking, not the legal
status of ranked firms. The ranking is useful for enforcement
only if it orders the loser-side firms for which bid-layer
follow-up is most valuable. The validation target is therefore
the one defined in Section~\ref{sec:cade_adjacency_labels}:
always-loser firms observed alongside direct CADE defendants in
adjudicated procurement-cartel environments. These cobidders are
adjudication-anchored exposure labels. They are not treated as
legal members of a cartel.

The validation exercise has to do more than report a high AUC.
The main empirical threat is exposure. Firms that participate
often, or that operate in the same procurement environments as
CADE defendants, have more opportunities to appear near those
anchors even under ordinary competition. Procurement-screen
designs therefore have to separate suspicious conduct from market
structure, bidder heterogeneity, and ordinary repeated
participation \citep{porter1993detection,bajari2003deciding}.
Network and large-scale procurement screens sharpen the same
concern \citep{conley2016detecting,kawai2022detecting}. A useful award-layer
screen must therefore concentrate cartel-adjacent loser-side
firms beyond what participation volume and opportunity-set exposure alone
would imply. A second threat is timing: an enforcement screen
should rank candidates before the target information is used.
A third threat is legal overreach: a loser-side screen should
not be read as a generic cartel-membership detector.

This section evaluates the ranking against those threats.
Subsection~\ref{sec:validation_baseline} reports the baseline
adjudication benchmarks. Subsection
\ref{sec:validation_exposure_timing} disciplines the result with
participation, exposure, leakage, and timing checks. Subsection
\ref{sec:validation_direct} uses direct CADE defendants as a
scope check. Subsection~\ref{sec:validation_summary} summarizes
what the validation establishes for the staged enforcement
architecture.

\subsection{Baseline Adjudication Benchmarks}
\label{sec:validation_baseline}

The baseline asks a simple question before adding stricter
discipline: do firms flagged by the award-layer rule appear
disproportionately among always-loser firms located near
adjudicated cartel anchors? The answer is yes. In the
conservative adjudication benchmark, we restrict the CADE
portfolio to cases resolved before the end of
$\valConservativeCutoff$. This leaves $\valConservativeCases$
cases, $\valConservativeFD$ BEC-active firm-defendants, and
$\valConservativeCobidders$ always-loser cobidders, of which
$\valConservativeFL$ are frequent losers.

% CO-AUTHOR EDIT [v21 Major 2]: disambiguate the 3.95% denominator.
% 3.95% is the FL->cobidder reach rate (108/2,735 FL firms), NOT the FL
% share among the 210 cobidders (which is 108/210=51.4%). Prose previously
% labeled 3.95% "share among conservative cobidders", a flat contradiction.
This benchmark uses early legal anchors and asks whether the
award-layer rank concentrates the associated loser-side
footprint. The enrichment quantity is the rate at which frequent
losers reach this conservative cobidder set: of the $\valFL$
frequent-loser firms, $\valConservativeFL$ co-bid with a
pre-$\valConservativeCutoff$ CADE defendant, a reach rate of
$\valCADEenrich$, against a participation-stratified placebo mean
of $\valCADEbase$. The excess ratio is
$\valMultThreeTwo\!\times$ with $p<0.001$. The corresponding ROC
AUC is $\valAUCprePost$ (95\% DeLong CI $\valAUCprePostCI$)
\citep{delong1988comparing}.
These numbers establish the baseline empirical fact: persistent
zero-win participation is not random with respect to
adjudication-anchored loser-side exposure.

The second benchmark adds timing discipline. The score is
computed from $2009$--$2016$ participation and evaluated against
$2017$--$2019$ cobidder labels. This is a temporal holdout, not a
real-time field deployment: it asks whether earlier award records
rank firms whose cartel adjacency is later observed in the
adjudication-anchored record. The firm-level AUC is
$\valAUCFLfirmTemp$ (95\% CI $\valAUCFLfirmTempCI$), below the
full-sample AUC $\valAUCFLfirm$ but well separated from random
ranking. Figure~\ref{fig:temporal_holdout_roc_submission} reports the
rolling-origin version of the same exercise.

The baseline results are deliberately read as ranking evidence.
They show that the award-layer statistic is informative about
where cartel-adjacent loser-side exposure later appears. The
next subsection tightens that reading by asking how much of the
signal remains after disciplining participation volume, market
exposure, leakage, and strict timing.

\begin{figure}[htbp]
\centering
\includegraphics[width=0.82\textwidth]{./output/figures/fig_temporal_holdout_roc}
\caption[Temporal-holdout discrimination]{Temporal-holdout ROC
curves. Scores are computed using participation observed before
the test year; the test year supplies the out-of-time cobidder
label. The exercise measures out-of-time discrimination against
later-adjudicated cobidder labels, not field deployment.}
\label{fig:temporal_holdout_roc_submission}
\end{figure}

\subsection{Exposure and Timing Discipline}
\label{sec:validation_exposure_timing}

The baseline result becomes useful for enforcement when the main
sources of mechanical exposure are made explicit. The first is
participation volume. In the conservative benchmark, always-loser
firms are grouped by participation quartiles, cobidder status is
permuted within quartiles $\valPermPilot$ times, and the share of
frequent losers that fall in the placebo cobidder set is
recomputed. This keeps the participation distribution fixed while
breaking the particular adjacency to CADE defendants. The observed
reach rate, $\valCADEenrich$, remains well above the placebo mean,
$\valCADEbase$. This audit
disciplines activity volume; it does not by itself make firms'
opportunities to meet CADE anchors comparable.

A formal version of the same volume placebo permutes the
frequent-loser label among always-losers while preserving the
marginal participation distribution. Across $B=2{,}000$ draws,
the sham AUC distribution has mean $\valShamMean$, standard
deviation $\valShamSD$, and 99th percentile $\valShamQNinetyNine$. The
observed FL14 AUC of $\valAUCFLfirm$ falls thirty-two standard
deviations above the sham mean, with permutation
$p$~$\valShamP$. Volume alone does not reproduce the observed
concentration. The price coefficient under the same sham
distribution does not reject the null---observed $\valShamPriceObs$ against
sham mean $\valShamPriceMean$, $p=\valShamPriceP$---which is consistent with
the scope-not-damages reading of the price evidence in
\S\ref{sec:price_scope_submission}.

The second source is market opportunity. Firms active in the
same procurement niches as CADE defendants have more chances to
appear near those anchors even when conduct is ordinary.
Exposure adjustment is therefore a required audit, not optional
robustness. The audit compares firms within more similar
opportunity sets---procurement cells (buyer, year, item-group) in
which a firm could have met a CADE anchor---and asks whether the
award-layer rank still concentrates adjudication-anchored
cobidders once the amount of opportunity is held fixed.
% CO-AUTHOR EDIT [v21 Major 1]: report the falsifiable number, not just prose.
Opportunity is itself highly predictive (AUC
$\valExpoExposureOnlyAUC$ from exposure alone), so the adjustment
is demanding. After it, the ranking is attenuated but does not
collapse: within opportunity strata the award-layer AUC is
$\valExpoWithinStratumAUC$, and adding the score to an
exposure-only model raises discrimination by $+\valExpoIncr$ AUC
(DeLong $p\approx\valExpoIncrP$). The unconditional $\valAUCFLfirm$
thus overstates the conduct-relevant signal, but loss intensity
still orders cobidders among comparably exposed firms. The audit
is not a legal-membership test and does not eliminate all
opportunity-set concerns. \ref{app:market_opportunity_submission}
reports the exact construction, the pre-committed rejection
region, and the retained support.

Timing is disciplined in two layers. The temporal holdout above
forms the participation score before the evaluation window and
tests later labels. A stricter version applies the same ex ante
logic to the screen definition itself: zero-win status, the
participation threshold, the frequent-loser flag, and the
continuous score are all formed using only pre-window records.
In that strict train-period audit, the threshold becomes
$\valThresholdTrain$ and the firm-level frequent-loser AUC is
$\valAUCStrictFirmFL$ (95\% CI $\valAUCStrictFirmFLCI$). The
continuous score gives AUC $\valAUCStrictFirmTC$ (95\% CI
$\valAUCStrictFirmTCCI$). The decline relative to the relaxed
holdout is informative: real-time screening has less information
than a full-window reconstruction, but the ranking remains
useful under a stricter field-like information set. It remains a
timing audit, not causal prediction.

Two structural facts make the timing audit harder to dismiss as
a same-firm artifact. First, firm persistence between the
$\valYearStart$--$2016$ and $2017$--$\valYearEnd$ windows is
$\valFirmPersist$: roughly nine in ten always-loser firms in the
test window are firms not seen in the training window. The
temporal holdout therefore evaluates a substantially fresh firm
population rather than the same firms in different years. Second,
when the training window is shortened further (scores frozen on
$2009$--$2015$, evaluated against cobidders linked to CADE
adjudications closed after $2019$---a strictly out-of-time
target), the firm-level AUC is $\valClfFifteenFL$ for the
frequent-loser flag and $\valClfFifteenPostTC$ for the
continuous score. The ranking remains informative when the
classifier is trained at one horizon, the cobidder label is
defined by adjudications at a later horizon, and the bulk of
test-period firms were not in the training set.
Table~\ref{tab:three_classifier_submission} reports the full
three-classifier battery.

\begin{table}[htbp]
\centering
\caption{Three-Classifier Timing Battery}
\label{tab:three_classifier_submission}
\footnotesize
\begin{adjustbox}{max width=\textwidth,center}
\begin{threeparttable}
\begin{tabular}{p{4.0cm}cccc}
\toprule
& \multicolumn{2}{c}{vs cobid\_all} & \multicolumn{2}{c}{vs cobid\_post2019} \\
\cmidrule(lr){2-3}\cmidrule(lr){4-5}
Classifier & FL flag & Continuous & FL flag & Continuous \\
\midrule
clf\_2015 (train 09--15) & $\valClfFifteenFL$ & $\valClfFifteenTC$
  & $\valClfFifteenPostFL$ & $\valClfFifteenPostTC$ \\
clf\_2017 (train 09--17) & $\valClfSeventeenFL$ & $\valClfSeventeenTC$
  & $\valClfSeventeenPostFL$ & $\valClfSeventeenPostTC$ \\
clf\_2019\_full (in-sample) & $\valAUCFLfirm$ & $\valAUClogtc$
  & --- & --- \\
\bottomrule
\end{tabular}
\begin{tablenotes}
\footnotesize
\item \textit{Notes:} AUC against two cobidder targets across three
classifiers. cobid\_all uses cobidders linked to any CADE adjudication
in the panel. cobid\_post2019 uses cobidders linked only to
adjudications closed after $2019$, a strictly out-of-time target
that cannot be in the clf\_2015 or clf\_2017 training data.
\end{tablenotes}
\end{threeparttable}
\end{adjustbox}
\end{table}

The leakage audit addresses a different concern: item-level
validation can mechanically reuse the same firm histories that
construct the score. We therefore treat the in-sample item-level
AUC, $\valLeakRefAUC$, as a mechanical reference point rather
than as the main validation estimate. Holding cobidder firms out
of score construction within cross-validation folds leaves AUC
$\valLeakCVAUC$ (95\% CI $\valLeakCVAUCCI$). The surviving
component is the signal relevant for triage: the award-layer
rank continues to order cartel-adjacent loser-side exposure when
the most direct reuse of firm labels in score construction is
weakened.

\subsection{Direct-Defendant Scope Check}
\label{sec:validation_direct}

The direct-defendant exercise aligns the validation target with
the screen's legal-economic role. Direct CADE defendants are the
legal anchors that locate adjudicated cartel environments. They
need not be persistent zero-win firms: in procurement cartels,
defendants may occupy winner-side, allocation, or rotating roles.
The award-layer statistic, by construction, ranks the loser side
of the procurement record. The empirical prediction is therefore
asymmetry: the rank should be informative for cobidders and much
less informative for direct defendants.

The data follow that prediction. Against direct defendants in
the broader BEC firm universe, AUC is $\valAUCdirectStd$, close
to random classification. This is the scope check doing its
work. If the statistic were a generic cartel-membership detector,
it should rank direct defendants as well as cobidders. It does
not, which is the expected result for a zero-win loser-side
ranking. Read together with the cobidder results, this pattern
supports the staged interpretation. Direct defendants anchor
cartel-affected environments; the award layer ranks the
loser-side exposure around those anchors; the bid layer then
supplies the conduct evidence for case-specific evaluation.

A complementary universe-anchored scope matrix systematically
varies the universe (always-loser firms vs full BEC firm
registry) and the positive class (cobidders vs direct defendants)
to map where the ranking targets succeed and where they fail. The
matrix produces a stronger result than a stand-alone null: raw
participation count on the full BEC universe against direct CADE
defendants returns AUC $\valAUCParticDefendants$, which is
\emph{below} random. A high-participation BEC firm is less
likely, not more likely, to be a direct CADE defendant under the
loser-side scoring logic. The score does not simply fail to
identify winner-heavy defendants; the relevant direction is
reversed, consistent with the role separation observed in D4
(\valDirectShareAL\% of direct defendants are always-losers; median win rate
$\valDirectMedWR$ versus $\valOthersMedWR$ for cobidders).
Table~\ref{tab:scope_matrix_submission} reports the full eight-row
scope matrix.

\begin{table}[htbp]
\centering
\caption{Universe-Anchored Stratum Scope Matrix}
\label{tab:scope_matrix_submission}
\footnotesize
\begin{adjustbox}{max width=\textwidth,center}
\begin{threeparttable}
\begin{tabular}{p{3.3cm}p{3.0cm}p{3.0cm}cp{4.0cm}}
\toprule
Score & Universe & Positive class & AUC & Reading \\
\midrule
Binary FL & Always-losers & Cobidders & $\valScopeRowOne$
  & Within-stratum prioritization \\
Continuous log\_tc & Always-losers & Cobidders & $\valScopeRowTwo$
  & Best within-stratum ranking \\
Binary FL & All BEC firms & Direct CADE & $\valScopeRowThree$
  & At chance; no membership detection \\
Participation count & All BEC firms & Direct CADE & $\valScopeRowFour$
  & \textbf{Below chance: actively repels winner-heavy} \\
Binary FL (frozen 09--16) & Always-losers as of 09--16 & Cobidders
  & $\valScopeRowFive$ & Weaker than pooled within-stratum \\
Continuous (train 09--16) & Items 17--19 & Cobidder items
  & $\valScopeRowSix$ & Some prospective generalization \\
Binary FL (frozen 09--16) & Items 17--19 & Cobidder items
  & $\valScopeRowSeven$ & Weak binary item-level holdout \\
Continuous (train 09--16) & Items 17--19 & Direct-defendant items
  & $\valScopeRowEight$ & Direct-defendant null survives timing \\
\bottomrule
\end{tabular}
\begin{tablenotes}
\footnotesize
\item \textit{Notes:} AUC across an eight-row matrix that varies the
score, the universe of firms or items being ranked, and the positive
class. Row~4 is below chance because the loser-side score
discriminates \emph{against} high-participation winner-heavy
defendants. The matrix structurally rules out a generic-detector
reading.
\end{tablenotes}
\end{threeparttable}
\end{adjustbox}
\end{table}

\subsection{Validation Summary}
\label{sec:validation_summary}

Table~\ref{tab:validation_benchmarks_submission} is a threat map
rather than a proof table. Each row corresponds to a threat
identified before the tests are read: legal anchoring,
participation volume, market exposure, timing, leakage, and
scope asymmetry. The resulting claim is staged: the award-layer
ranking supports a priority queue for the bid-layer stage, where
tender-level conduct can be examined with richer evidence.

\begin{table}[htbp]
\centering
\caption{Validation Architecture for the Award-Layer Ranking}
\label{tab:validation_benchmarks_submission}
\footnotesize
\begin{adjustbox}{max width=\textwidth,center}
\begin{threeparttable}
\begin{tabular}{p{3.5cm}p{5.3cm}p{4.3cm}}
\toprule
Benchmark & Question & Result \\
\midrule
Legal anchoring & Are frequent losers overrepresented among pre-$\valConservativeCutoff$ adjudication-anchored cobidders? & $\valMultThreeTwo\!\times$ excess; AUC $\valAUCprePost$ \\
Participation-volume discipline & Does concentration survive volume-preserving placebo assignment? & Placebo mean $\valCADEbase$ versus observed $\valCADEenrich$ \\
Exposure discipline & Does the ranking remain informative under more comparable opportunity sets for meeting CADE anchors? & Within-opportunity AUC $\valExpoWithinStratumAUC$; $+\valExpoIncr$ over exposure-only ($p\approx\valExpoIncrP$) \\
Temporal holdout & Does pre-window participation rank later cobidder labels? & AUC $\valAUCFLfirmTemp$ \\
Leakage audit & Does signal survive when cobidders are held out of score construction? & Out-of-fold AUC $\valLeakCVAUC$ \\
Direct-defendant scope check & Does the ranking preserve the distinction between legal anchors and loser-side exposure? & Direct-defendant AUC $\valAUCdirectStd$ \\
\bottomrule
\end{tabular}
\begin{tablenotes}
\footnotesize
\item \textit{Notes:} Rows address distinct threats to
interpreting the ranking as forensic priority. The benchmarks
validate a legally limited ranking, not cartel membership,
liability, or proof.
\end{tablenotes}
\end{threeparttable}
\end{adjustbox}
\end{table}

Taken together, the benchmarks support the screen at the stage
for which it is designed. Persistent zero-win participation
orders adjudication-anchored loser-side exposure well enough to
allocate forensic attention, not to validate the legal status of
ranked firms. The next step is economic interpretation: among
firms already selected by the frequent-loser rule, do cobidders
display a distinct loser-side profile rather than a mechanical
collection of high-volume losers?
